\HMTNarrativeHeading{One unity, two realizations}

The double projection organizes this article from its first construction.
One and the same state preserves the data that allow a value to be refined
and the data that allow an incidence to be determined. The first realization
assembles compatible cylinders until its numerical reading is uniquely
determined; the second transports boundary relations. The thesis examines
their common origin and the operations that maintain their correspondence
throughout refinement. Number and geometrical configuration thereby acquire
an explicit genealogy.

Unity is constituted through positions, operations, and transport relations.
APP arranges the two nonadic directions on a common support and evaluates its
additive and multiplicative sheets. Each evaluation preserves the remainder
and its quotient. TRIT distinguishes orientation and the local regime; the TPK
composes these data with displacement, carry, boundary, and prolongation.
The joint discrete structure of the continuum resides in this enriched
construction. Its subsequent realizations receive operations that have already
been determined.

The nonadic connection has a constitutive function here. One turn composes
the transitions of the complete state. The observable phase returns and the
register incorporates the path traversed. This distinction makes it possible
to recognize a regularity while the history increases in depth. The continuity
of refinement preserves the prefixes that make the two realizations
compatible. Each new level of precision belongs to the same history from
which the previous levels originate.

The meaning of evolutionary conservation becomes apparent at this point.
A normalized quantity may remain unchanged while its distribution,
resolution, or representation changes. Recovering the state requires
identifying the data that accompany that change. Arithmetic preserves
quotients; transport preserves orientation and turns; incidence preserves
relations; the operator archive preserves the complements needed to invert
the reading. The article develops the maps that articulate these functions.

The holographic extreme and mean ratio expresses this question of unity,
proportion, and transformation. Its formulation includes dimensional
completion and reversible refinement, and subsequently extends to balances
of states and observables. The coupling constant occupies a definite position
within this development: the generated register links the arithmetic reading,
incidence, and operator recovery. The proofs that follow construct the
antecedents of each of these functions.
